\section{Introduction}
Every production system fails: caches crash, databases slow down, networks drop packets and
processes die. Operations teams respond with runbooks, typically of the form ``if $X$ is broken,
restart $Y$'', and increasingly with automation and AI assistants that act on those runbooks.
Three gaps make this fragile. First, runbooks are rarely \emph{measured}: whether ``restart it''
actually restores service, and how fast, is assumed rather than known. Second, automation can
amplify incidents: restart loops and actions on the wrong component are classic causes of
outages, so autonomous remediation needs guardrails. Third, release pipelines test functional
correctness on a healthy system; they rarely test that the system still \emph{recovers} when a
dependency fails.

Chaos engineering~\cite{basiri2016chaos,principlesofchaos} addresses part of the problem by
injecting faults on purpose, and Site Reliability Engineering (SRE) supplies the vocabulary to
judge the outcome: service level indicators (SLIs), objectives (SLOs), error budgets and burn-rate
alerts~\cite{beyer2016sre,beyer2018workbook}. Research on AIOps applies learning to failure
management~\cite{dang2019aiops,notaro2021survey}, and recent work uses large language models
(LLMs) to recommend root causes and mitigations for cloud
incidents~\cite{ahmed2023recommending,chen2024rcacopilot,roy2024agents}. What is missing is a
place where these ideas meet \emph{under measurement}: where competing remediation strategies,
including learned and LLM-based ones, act on real faults behind an explicit safety layer, where
every recovery is timed per request, and where the evidence gates software releases.

This paper presents \sys{}, an open-source resilience lab built from standard DevOps components
(Docker Compose, Ansible, Jenkins, Prometheus, Alertmanager, Grafana) around a three-tier web
service. It makes six contributions:
\begin{enumerate}
\item A \textbf{measurement harness} that injects seven real faults under open-model load and
  defines recovery precisely at request granularity (Section~\ref{sec:measure}).
\item A \textbf{safety shield} that checks every proposed action, from any policy, against a
  declarative policy (environment, allowlist, no repeats, action budget, approval), backed by a
  second allowlist in the Ansible playbooks, plus human approval in on-call mode
  (Section~\ref{sec:safety}).
\item A \textbf{comparison of five remediation policies} (a conventional runbook, expert rules, a
  LinUCB contextual bandit~\cite{li2010linucb}, gpt-4o-mini, and a tool-using gpt-4o-mini agent)
  on training and held-out faults, with a learning curve for the bandit (Sections~\ref{sec:policies}
  and~\ref{sec:results}).
\item A \textbf{2$\times$2 ablation} that separates the effect of giving an LLM investigation tools
  from the effect of generic SRE guidance in its prompt.
\item An \textbf{SLO-driven on-call agent} that is paged by multi-window burn-rate alerts,
  confirms the page with its own probe traffic, remediates through the shield, escalates instead of
  looping, and writes AI post-mortems, evaluated on unannounced faults (Section~\ref{sec:oncall}).
\item A \textbf{CI/CD pipeline} in which resilience, alert-rule, security and on-call checks are
  release gates; a one-character bug that passes unit and smoke tests is blocked by the resilience
  gate (Section~\ref{sec:cicd}).
\end{enumerate}
